% Historical figure blocks; not included by the active manuscript.

\begin{figure}[!htbp]
    \centering
    \includegraphics[width=\linewidth]{arxiv_pics/fig14_memory_power_day5_v5_pending122.pdf}
    \caption{\vfourred{Day-5 physical power spectra for T2m, Z500, Q700, and six-hour precipitation. ERA5 is compared with raw and filtered GraphCast, M1, and reset/reset. Each panel retains its own physical squared units. Powers are pooled over dates; M1 and reset/reset are shown without amplitude rescaling.}}
    \label{fig:memory_power_day5}
\end{figure}

\begin{figure}[!htbp]
    \centering
    \includegraphics[width=\linewidth]{arxiv_pics/fig15_memory_power_day10_v5_pending122.pdf}
    \ifvfivechanges
    \caption*{\color{vfiveDeleted}\sout{Day-10 physical power spectra for the same forecasts, fields, and conventions as Figure~\mbox{\ref{fig:memory_power_day5}}.}}
    \fi
    \caption{\vfourred{Provisional day-10 physical power spectra of (a) $2$\,m temperature (T2m), (b) geopotential at 500\,hPa (Z500), (c) specific humidity at 700\,hPa (Q700), and (d) six-hour accumulated precipitation (TP6h). Curves compare ERA5, raw and filtered GraphCast, M1, and reset/reset; each panel shows its variable's physical power units on the vertical axis.}}
    \label{fig:memory_power_day10}
\end{figure}
